\documentclass[10pt]{article}
\usepackage[utf8]{inputenc}\usepackage[T1]{fontenc}
\usepackage[spanish,es-nodecimaldot]{babel}\usepackage{amsmath}
\usepackage[paperwidth=485pt,paperheight=1800pt,margin=1pt]{geometry}
\setlength{\parindent}{0pt}\setlength{\parskip}{5pt}\pagestyle{empty}
\begin{document}
La variable original está en kelvin y se convierte a grados Celsius mediante 
\begin{equation}
T[{}^\circ\mathrm{C}]=T[\mathrm{K}]-273{,}15\tag{3.1}
\end{equation}
 Para reproducir la descarga se seleccionan las medias mensuales (monthly averaged), la variable 2m temperature, todos los meses de 1981 a 2026 y formato NetCDF. La petición original no conservó el área; para repetirla se recomienda norte 4,85$^\circ$, oeste 75,95$^\circ$O, sur 4,05$^\circ$N y este 75,35$^\circ$O. El archivo disponible contiene 548 meses, enero de 1981--agosto de 2026, con una malla de 0,1$^\circ$(aproximadamente 11 km; resolución nativa cercana a 9 km).
\newpage
El archivo ya contiene medias mensuales de medias diarias: no se suman ni se vuelve a calcular una media temporal. Si se parte de datos diarios completos, la media mensual se calcula como
\begin{equation}
T_m=\frac{1}{N_m}\sum_{d=1}^{N_m}T_d\tag{3.2}
\end{equation}
Se utilizan todos los 28, 29, 30 o 31 días del mes. Los 548 meses de la serie de cuenca están válidos. Para representar la cuenca se ponderan 36 celdas según el área de intersección con el polígono:
\begin{equation}
T_m^B=\sum_i w_iT_{i,m},\qquad w_i=\frac{A(B\cap C_i)}{\sum_i A(B\cap C_i)}\tag{3.3}
\end{equation}
Las celdas cubren el 100 \% del polígono de cuenca.
\newpage
Para cada mes calendario j, la climatología se calculó sobre todos los años disponibles del periodo de referencia. Se usó la desviación estándar muestral (n -1) y las definiciones 
\begin{equation}
a_t=X_t-\mu_{j(t)},\qquad z_t=\frac{a_t}{s_{j(t)}}\tag{3.4}
\end{equation}
 El mismo periodo de referencia se aplica a los doce meses, de modo que las anomalías comparan cada observación con su estación del año y no eliminan una tendencia temporal mediante una climatología móvil.
\newpage
Secuencia cronológica completa. Para evitar atribuir al tiempo el ciclo anual, el modelo principal incluye efectos fijos para los doce meses calendario y una pendiente temporal común: 
\begin{equation}
Y_t=\alpha_{j(t)}+\beta t+\epsilon_t\tag{3.5}
\end{equation}
 La misma especificación se aplicó a la serie original, la anomalía y la anomalía estandarizada. Como contraste, se muestra también la pendiente OLS ingenua sin control estacional; en la precipitación original puede mezclar ciclo anual y evolución temporal.
\newpage
1. Tendencia lineal por OLS (paramétrica). Se ajustó 
\begin{equation}
Y_t=\beta_0+\beta_1t+\epsilon_t\tag{3.6}
\end{equation}
 y sus equivalentes en at y zt. En la serie original se contrastó ese OLS simple con una regresión que añade efectos fijos de mes calendario y una pendiente temporal común. El modelo estacional controla el ciclo anual; sus resultados son la comparación principal. Para los intervalos y pruebas se usaron errores HAC Newey--West con rezago 12 meses en la secuencia completa y 1 año en las subseries. El diagnóstico de residuos reporta autocorrelación de rezagos 1 y 12, Ljung-- Box hasta 12 meses, cambio temporal de varianza con efectos de mes y razón entre desviaciones estándar de las mitades del registro.
\newpage
Transformada, ventana y normalización. En cada tramo continuo de N meses se calcula la transformada discreta:
\begin{equation}
X_k=\sum_{t=0}^{N-1}x_t\exp\!\left(-\frac{2\pi\mathrm{i}kt}{N}\right)\tag{4.1}
\end{equation}
Se reportan las frecuencias positivas, los periodos equivalentes y el espaciamiento entre frecuencias:
\begin{equation}
f_k=\frac{k}{N\Delta t},\qquad T_k=\frac{1}{f_k},\qquad\Delta f=\frac{1}{N\Delta t}\tag{4.2}
\end{equation}
Se usa el periodograma unilateral con ventana rectangular (sin taper), cuya potencia de base se normaliza como
\begin{equation}
\mathcal{P}_k=\frac{|X_k|^2}{N^2}\tag{4.3}
\end{equation}
Esta potencia se duplica en las frecuencias positivas, excepto en Nyquist cuando N es par. La suma de potencias reproduce la varianza poblacional de la señal centrada y las unidades son el cuadrado de la variable. No se rellena con ceros ni se interpreta la frecuencia cero como un periodo finito.
\newpage
Convención y lectura. Se presentan periodogramas unilaterales de densidad espectral. Para una frecuencia de muestreo de una muestra por mes, la densidad de base se expresa como
\begin{equation}
S(f_k)=\frac{|X_k|^2}{f_s\sum_t w_t^2},\qquad f_s=1\ \mathrm{muestra/mes}\tag{4.4}
\end{equation}
Se duplican los bins positivos, salvo Nyquist. La densidad tiene unidades de variable al cuadrado por mes. La potencia integrada por bin y su normalización porcentual son
\begin{equation}
\mathcal{P}_k=S(f_k)\Delta f,\qquad p_k=100\,\frac{\mathcal{P}_k}{\sum_{j>0}\mathcal{P}_j}\tag{4.5}
\end{equation}
La potencia por bin tiene unidades de variable al cuadrado. Las curvas normalizadas permiten comparar formas, no variabilidad absoluta entre variables. La varianza original y su unidad se conservan en los archivos de resultados.
\newpage
Las medias mensuales no se suman. SLP y presión superficial vienen en las unidades declaradas en el NetCDF: se dividen por 100 solo si están en Pa; millibars equivale a hPa. NCEP hgt ya es altura geopotencial en metros: no se vuelve a dividir por g. En ERA5, geopotential es $\Phi$ en m$^2$/s$^2$ y se convertiría a 
\begin{equation}
Z=\frac{\Phi}{g},\qquad g=9{,}80665\ \mathrm{m/s^2}\tag{5.1}
\end{equation}
La altura resultante se expresa en metros.  Se conservan los faltantes oceánicos/terrestres de ERSST sin interpolarlos. Las longitudes se normalizan a [-180$^\circ$,180$^\circ$) sin duplicar la costura; cada campo mantiene su malla nativa.
\newpage
La referencia comprende las mismas 285 fechas completas de cuenca en 1998--2022. Para cada variable se resta la climatología de su mes calendario; los campos globales y las respuestas conservan sus unidades antes de calcular el coeficiente adimensional. La base tiene rezago cero. Para Q se añade un rezago de un mes como exploración de la persistencia climática y de una respuesta hidrológica con almacenamiento; no se interpreta como tiempo de concentración. Con la convención 
\begin{equation}
r_j(\lambda,\varphi;\ell)=\operatorname{corr}\!\left[a_X(t),a_Y(\lambda,\varphi,t-\ell)\right],\qquad j(t)=j\tag{5.2}
\end{equation}
 ell>0 significa que el campo climático antecede al caudal. Enero de 1998 se empareja con diciembre de 1997 para ell=1. El desplazamiento se hace en el calendario completo; no se usan valores futuros ni se comprimen los huecos.
\newpage
La dependencia relevante se estima entre años consecutivos del mismo mes. No se trata enero y febrero como observaciones consecutivas ni se conectan artificialmente años separados por huecos. Por celda, rhoX y rhoY son las autocorrelaciones anuales de orden uno de los mismos pares; se requieren al menos diez pares anuales adyacentes. Se adopta la aproximación AR(1) 
\begin{equation}
n_{\mathrm{ef}}=n\,\frac{1-\rho_X\rho_Y}{1+\rho_X\rho_Y},\qquad 3\leq n_{\mathrm{ef}}\leq n\tag{5.3}
\end{equation}
Limitada al intervalo indicado y  basada en el tamaño efectivo para correlaciones de Bretherton et al. (1999). El límite superior evita aumentar la información por autocorrelación estimada negativa. La prueba bilateral usa 
\begin{equation}
t=|r|\sqrt{\frac{\nu}{1-r^2}}\tag{5.4}
\end{equation}
Los grados de libertad utilizados son
\begin{equation}
\nu=\begin{cases}n_{\mathrm{ef}}-2,&\text{para anomalías},\\n_{\mathrm{ef}}-3,&\text{al controlar el año}.\end{cases}\tag{5.5}
\end{equation}
 Con un grado de libertad o menos no se emite p. Los p son aproximados: n corto, no normalidad, incertidumbre de rho y memoria de orden superior pueden afectar su calibración; no son una garantía exacta de cobertura o significancia.
\newpage
\begin{equation}
e_t=P_{\mathrm{IMERG},t}-L_t\tag{2.1}
\end{equation}
\newpage
\begin{equation}
\rho(k)=\operatorname{corr}\!\left[P(t-k),Q(t)\right],\qquad k\geq 0\tag{2.2}
\end{equation}
\newpage
\begin{equation}
R_m=\frac{86{,}4\,d_m\,Q_m}{2797{,}19}\quad[\mathrm{mm/mes}]\tag{2.3}
\end{equation}
\newpage
\begin{equation}
X'(t)=X(t)-\overline{X}_{j(t)}\tag{2.4}
\end{equation}
\newpage
\begin{equation}
Q_t=a+bP_t\tag{2.5}
\end{equation}
\newpage
\begin{equation}
Q_t=a+bP_{t-1}\tag{2.6}
\end{equation}
\newpage
\begin{equation}
Q_t=a+bP_t+cP_{t-1}\tag{2.7}
\end{equation}
\newpage
\begin{equation}
Q_t=a+b\sqrt{P_t}\tag{2.8}
\end{equation}
\newpage
\begin{equation}
\widehat{Q}_t=\max\left(0,-87.2692+0.465439P_{IMERG,t}+0.433065P_{IMERG,t-1}\right).\tag{2.9}
\end{equation}
\newpage
\begin{equation}
\widehat{Q}_t=\max\left(0,-35.8526+0.422155P_{CHIRPS,t}+0.434939P_{CHIRPS,t-1}\right).\tag{2.10}
\end{equation}
\newpage
\begin{equation}
\widehat{L}_t=\max\left(0,-30.7549+0.596818P_{IMERG,t}\right).\tag{2.11}
\end{equation}
\newpage
\begin{equation}
\widehat{L}_t=\max\left(0,-0.9713+0.504221P_{CHIRPS,t}\right).\tag{2.12}
\end{equation}
\newpage
\begin{equation}
\widehat{Q}_t=\max\left(0,55.3322+0.654768L_t\right).\tag{2.13}
\end{equation}
\newpage
\begin{equation}
\widehat{R}_m=\frac{86{,}4\,d_m\,\widehat{Q}_m}{2797{,}19}\tag{2.14}
\end{equation}
\newpage
\begin{equation}
\widehat{Q}_{\mathrm{IMERG},t}=\max\!\left(0,-87{,}269198+0{,}465439P_t+0{,}433065P_{t-1}\right)\tag{2.15}
\end{equation}
\newpage
\begin{equation}
\widehat{Q}_{\mathrm{CHIRPS},t}=\max\!\left(0,-35{,}852625+0{,}422155P_t+0{,}434939P_{t-1}\right)\tag{2.16}
\end{equation}
\newpage
\begin{equation}
\mathrm{NSE}=1-\frac{\mathrm{SSE}}{\mathrm{SST}}\tag{2.17}
\end{equation}
\newpage
\begin{equation}
T[{}^\circ\mathrm{C}]=T[\mathrm{K}]-273{,}15\tag{3.1}
\end{equation}
\newpage
\begin{equation}
T_m=\frac{1}{N_m}\sum_{d=1}^{N_m}T_d\tag{3.2}
\end{equation}
\newpage
\begin{equation}
T_m^B=\sum_i w_iT_{i,m},\qquad w_i=\frac{A(B\cap C_i)}{\sum_i A(B\cap C_i)}\tag{3.3}
\end{equation}
\newpage
\begin{equation}
a_t=X_t-\mu_{j(t)},\qquad z_t=\frac{a_t}{s_{j(t)}}\tag{3.4}
\end{equation}
\newpage
\begin{equation}
Y_t=\alpha_{j(t)}+\beta t+\epsilon_t\tag{3.5}
\end{equation}
\newpage
\begin{equation}
Y_t=\beta_0+\beta_1t+\epsilon_t\tag{3.6}
\end{equation}
\newpage
\begin{equation}
X_k=\sum_{t=0}^{N-1}x_t\exp\!\left(-\frac{2\pi\mathrm{i}kt}{N}\right)\tag{4.1}
\end{equation}
\newpage
\begin{equation}
f_k=\frac{k}{N\Delta t},\qquad T_k=\frac{1}{f_k},\qquad\Delta f=\frac{1}{N\Delta t}\tag{4.2}
\end{equation}
\newpage
\begin{equation}
\mathcal{P}_k=\frac{|X_k|^2}{N^2}\tag{4.3}
\end{equation}
\newpage
\begin{equation}
S(f_k)=\frac{|X_k|^2}{f_s\sum_t w_t^2},\qquad f_s=1\ \mathrm{muestra/mes}\tag{4.4}
\end{equation}
\newpage
\begin{equation}
\mathcal{P}_k=S(f_k)\Delta f,\qquad p_k=100\,\frac{\mathcal{P}_k}{\sum_{j>0}\mathcal{P}_j}\tag{4.5}
\end{equation}
\newpage
\begin{equation}
Z=\frac{\Phi}{g},\qquad g=9{,}80665\ \mathrm{m/s^2}\tag{5.1}
\end{equation}
\newpage
\begin{equation}
r_j(\lambda,\varphi;\ell)=\operatorname{corr}\!\left[a_X(t),a_Y(\lambda,\varphi,t-\ell)\right],\qquad j(t)=j\tag{5.2}
\end{equation}
\newpage
\begin{equation}
n_{\mathrm{ef}}=n\,\frac{1-\rho_X\rho_Y}{1+\rho_X\rho_Y},\qquad 3\leq n_{\mathrm{ef}}\leq n\tag{5.3}
\end{equation}
\newpage
\begin{equation}
t=|r|\sqrt{\frac{\nu}{1-r^2}}\tag{5.4}
\end{equation}
\newpage
\begin{equation}
\nu=\begin{cases}n_{\mathrm{ef}}-2,&\text{para anomalías},\\n_{\mathrm{ef}}-3,&\text{al controlar el año}.\end{cases}\tag{5.5}
\end{equation}
\end{document}